\documentclass[a4paper]{article}
% Español (México) con XeLaTeX: fontspec para fuentes Unicode, polyglossia
% para la división silábica y los nombres del documento en la variante mexicana.
\usepackage{fontspec}
\usepackage{polyglossia}
\setdefaultlanguage[variant=mexican]{spanish}
% Los nombres chinos van como «pinyin (original)»: xeCJK da a los caracteres
% Han su propia fuente; sin ella XeLaTeX los omite con un aviso.
% FandolSong no cubre algunos caracteres de nombres (U+73FA); WenQuanYi Zen Hei sí.
\usepackage[AutoFallBack=true]{xeCJK}
\setCJKmainfont{FandolSong-Regular.otf}
\setCJKfallbackfamilyfont{\CJKrmdefault}{WenQuanYi Zen Hei}
% polyglossia no traduce los nombres de listings.
\AtBeginDocument{\providecommand{\lstlistingname}{}\renewcommand{\lstlistingname}{Listado}%
  \providecommand{\lstlistlistingname}{}\renewcommand{\lstlistlistingname}{Índice de listados}}
\usepackage{amsmath, amssymb}
\usepackage{graphicx}
\usepackage[margin=2.35cm]{geometry}
\usepackage[most]{tcolorbox}
\usepackage{booktabs}
\usepackage{tabularx}
\usepackage{longtable}
\usepackage{array}
\usepackage{float}
\usepackage{xcolor}
\usepackage{hyperref}
\usepackage{fancyhdr}
\hypersetup{colorlinks=true, linkcolor=blue!55!black, urlcolor=blue!60!black}
\graphicspath{{./}{figures/}}
\setlength{\parskip}{0.55em}
\setlength{\parindent}{2em}
\setlength{\emergencystretch}{3em}
\renewcommand{\arraystretch}{1.18}
\newtcolorbox{knowledgebox}[1]{enhanced, colback=blue!5!white, colframe=blue!70!black, colbacktitle=blue!70!black, coltitle=white, fonttitle=\bfseries, title={#1}, sharp corners, breakable}
\newtcolorbox{importantbox}[1]{enhanced, colback=yellow!10!white, colframe=yellow!75!black, colbacktitle=yellow!75!black, coltitle=black, fonttitle=\bfseries, title={#1}, sharp corners, breakable}
\newtcolorbox{warningbox}[1]{enhanced, colback=red!5!white, colframe=red!70!black, colbacktitle=red!70!black, coltitle=white, fonttitle=\bfseries, title={#1}, sharp corners, breakable}
\pagestyle{fancy}
\fancyhf{}
\fancyfoot[C]{\thepage}
\fancyfoot[R]{\footnotesize\color{gray}\href{https://github.com/hqhq1025/ai-course-notes}{AI Course Notes}}
\renewcommand{\headrulewidth}{0pt}
\renewcommand{\footrulewidth}{0pt}
\newcommand{\videourl}{https://www.youtube.com/watch?v=u1Lzp-7Ybn8}
\newcommand{\repourl}{https://github.com/hqhq1025/ai-course-notes}

\begin{document}
\begin{titlepage}
\centering
{\Large Notas de revisión técnica\par}
\vspace{1.0cm}
{\huge\bfseries Informe trimestral global de grandes modelos, episodio 9: Coding, Agent y el modelo como OS\par}
\vspace{0.6cm}
{\Large Zhang Xiaojun conversa con Guang Mi (广密)\par}
\vspace{0.4cm}
{\large Elaborado a partir del video público de Zhang Xiaojun Podcast, una transcripción hecha en GPU y diagramas conceptuales propios\par}
\vspace{0.3cm}
{\large 2026-04-15\par}
\vspace{0.9cm}
\includegraphics[width=0.88\textwidth,height=0.42\textheight,keepaspectratio]{cover.jpg}\par
\vfill
\begin{tcolorbox}[width=0.92\textwidth, colback=black!2!white, colframe=black!55, sharp corners]
\textbf{Título del video}: 136. Informe trimestral global de grandes modelos, episodio 9: conversación con Guang Mi. Coding es el segundo acto de la AGI, la verdad sobre las tres grandes de Silicon Valley y el modelo se está convirtiendo en el OS de nueva generación\par
\textbf{Canal del video}: Zhang Xiaojun Podcast\par
\textbf{Duración del video}: 01:22:40\par
\textbf{Enlace del video}: \href{\videourl}{\nolinkurl{\videourl}}\par
\textbf{Nota sobre el material}: YouTube no ofrece subtítulos disponibles; el texto se elaboró a partir de la transcripción con faster-whisper large-v3 en CUDA, los capítulos del video, la portada y figuras didácticas propias. Las tomas fijas de la entrevista no se repiten en el cuerpo del texto.\par
\textbf{Página del proyecto}: \href{\repourl}{\nolinkurl{\repourl}}\par
\end{tcolorbox}
\end{titlepage}

\tableofcontents
\newpage

\section{Introducción: por qué el tono de este informe trimestral es ambivalente}

El tono central de este episodio del «informe trimestral global de modelos grandes» tiene dos caras. Por un lado, Coding y Agentic workflow están llevando a la IA de los chatbots a sistemas capaces de realizar trabajo, y tanto la capacidad de los modelos como la eficiencia de la investigación se aceleran. Por otro lado, esta aceleración está empujando el trabajo de oficina, la organización de la investigación y desarrollo y la distribución social hacia un periodo de deflación y desempleo. El juicio central de Guangmi (广密) es que Coding es el nuevo acelerador de la IA: un modelo líder en Coding es como una GPU de vanguardia y se convertirá en un amplificador decisivo del avance hacia la AGI.

\begin{importantbox}{Tesis central de este episodio}
Si el Chatbot es el primer acto y Coding/Agent es el segundo, en la siguiente etapa las empresas de modelos no compiten solo por «quién conversa mejor», sino por quién logra conectar el modelo con el código, las herramientas, los entornos de tareas y los flujos de trabajo reales, y convertir esa capacidad en una nueva generación de sistema operativo.
\end{importantbox}

\begin{knowledgebox}{Nota sobre la estrategia visual}
Este video muestra la imagen fija de una entrevista o pódcast de audio, sin slides didácticas, pizarrón ni demostraciones de producto. Conforme al estándar de pódcast de este repositorio, el cuerpo del texto no repite fotogramas de los participantes; la portada sirve para identificar la fuente y el cuerpo del texto expone la estructura del panorama mediante diagramas conceptuales y tablas.
\end{knowledgebox}

\subsection{Resumen de la sección}

Este episodio no es una crónica de noticias, sino una actualización trimestral del marco de análisis: Coding se convierte en acelerador de la AGI, se reordenan los periodos de competencia entre las empresas de modelos, el modelo podría convertirse en la nueva generación de OS y el impacto social empieza a pasar de la «eficiencia de las aplicaciones» a la «estructura de los puestos de trabajo».
\section{Coding es el segundo acto de la AGI}

La afirmación más contundente de la entrevista es que Coding lleva a la IA del primer acto, el del Chatbot, al segundo acto, el del Agent. El núcleo del Chatbot son las preguntas y respuestas, la búsqueda, el resumen y la conversación; el núcleo de Coding es convertir una intención expresada en lenguaje natural en una solución ejecutable, de modo que el modelo modifique código, ejecute tareas, invoque herramientas, procese datos y construya sistemas directamente. Con ello, la IA pasa de «describir el mundo» a «transformar el mundo digital».

\begin{figure}[H]
\centering
\includegraphics[width=0.94\textwidth]{coding-second-act.png}
\caption{Coding es el segundo acto de la AGI: de conversar a trabajar, y de ahí al modelo como OS.}
\end{figure}

\begin{knowledgebox}{Lectura de la figura: por qué Coding es el punto de inflexión}
En la figura, el primer acto es el Chatbot, cuyo valor proviene de la compresión de información; el segundo acto es Coding/Agent, cuyo valor proviene de completar trabajo; el tercer acto es Model as OS, en el que el modelo se convierte en la capa de coordinación entre la intención del usuario y la ejecución de las aplicaciones. Lo decisivo de Coding no es escribir código en sí, sino que el código puede expresar la mayoría de las soluciones del mundo digital.
\end{knowledgebox}

\subsection{«El lenguaje es el mundo, el código es la solución»}

En la entrevista hay una frase muy adecuada como título de esta sección: el lenguaje natural es una descripción del mundo, y el código es una descripción de la solution. El lenguaje natural puede expresar intenciones; el código puede expresar rutas de ejecución. Si un modelo puede escribir código de manera confiable, no solo puede responder preguntas, sino también manipular archivos, invocar API, procesar datos, automatizar procesos y construir productos. Por eso los modelos de Coding se comparan con la GPU: no son una aplicación aislada, sino un acelerador.

\begin{importantbox}{Coding como acelerador}
La GPU acelera el entrenamiento de modelos; los modelos de Coding aceleran el ciclo de investigación e ingeniería de las personas y de la IA. Que una idea pase de tardar dos o tres semanas en funcionar a tardar dos o tres días significa que el volumen de experimentos, la iteración del producto y la retroalimentación de la investigación se comprimen.
\end{importantbox}

\subsection{Los investigadores ya no escriben el código ellos mismos}

Guangmi (广密) mencionó que los investigadores de los laboratorios de vanguardia y los programadores de alto nivel ya escriben mucho menos código por su cuenta y han pasado a un esquema de «la IA escribe, la persona revisa». Si la IA alcanza una capacidad de implementación al nivel de un CTO o de un arquitecto en jefe, el trabajo humano pasará de escribir la implementación a fijar la dirección, revisar el diseño, verificar los resultados y asumir la responsabilidad. Esto se corresponde con la puesta en producción de Agent del EP139 y con el sistema de posentrenamiento del EP138.

\begin{warningbox}{No hay que entender «las personas no escriben código» como «las personas no importan»}
Que las personas no escriban grandes cantidades de código de implementación no significa que se retiren de la ingeniería. Las capacidades verdaderamente escasas se desplazarán hacia definir problemas, identificar errores, diseñar experimentos, organizar sistemas y juzgar el valor. La capacidad de revisión, la capacidad de arquitectura y la capacidad de definir tareas se vuelven, de hecho, más importantes.
\end{warningbox}

\subsection{Resumen de la sección}

Coding es el segundo acto porque convierte a la IA de un sistema que produce lenguaje en un sistema que actúa. No solo aumenta la eficiencia de los programadores, sino que también acelera la propia investigación en IA y se convierte en el amplificador central en el camino hacia la AGI.
\section{Percepción desde Silicon Valley: la ventana de progreso de los modelos se reordena}

La percepción de Guangmi (广密) es que, en el último trimestre, el nivel de los modelos avanzó más que en todo 2025; saltos como el de Opus 4.5 a 4.6 llevaron a los modelos de las preguntas y respuestas tipo chat a un modo verdaderamente agentic. La entrevista sitúa a las «tres grandes» de Silicon Valley en una ventana dinámica: OpenAI, Anthropic y Gemini tienen cada una su propia ventana de oportunidad, y la fórmula del éxito de hoy puede convertirse en el veneno de la siguiente etapa.

\begin{figure}[H]
\centering
\includegraphics[width=0.96\textwidth]{frontier-lab-map.png}
\caption{Mapa de posiciones de las empresas de modelos de Silicon Valley según las opiniones de la entrevista.}
\end{figure}

\begin{knowledgebox}{Lectura de la figura: no es una clasificación objetiva, sino un mapa de dependencia de la trayectoria}
La figura organiza las opiniones de la entrevista: Anthropic tiene una buena ventana en Coding/Agent; el éxito de OpenAI en ToC puede frenar su estrategia de Coding; Gemini tiene recursos y un ecosistema sólidos, pero le falta impulso; Meta se considera un retador; xAI se ve afectada por los vaivenes estratégicos y por problemas de equipo. El lector debe fijarse en la dependencia de la trayectoria, no tomarla como una clasificación fija.
\end{knowledgebox}

\subsection{Anthropic: detalles técnicos y ADN cultural}

En la entrevista se considera que Anthropic aprovechó la ventana de Coding/Agent. No tenía todo claro desde el day 1, pero sus fundadores revisan los datos de forma hands-on, dan importancia a los detalles técnicos y aplican entrevistas culturales estrictas; todo ello le facilita organizar recursos cuando surge un nuevo paradigma. Lo importante aquí es que la cultura no es un eslogan: influye en si el equipo valora cierto tipo de datos, de productos y de capacidades.

\subsection{OpenAI: las victorias pasadas pueden convertirse en veneno}

El éxito de OpenAI en ToC con ChatGPT podría, paradójicamente, llevar a la organización a concentrarse durante mucho tiempo en la puerta de entrada conversacional y en los productos de consumo, y a descuidar Coding. El juicio expresado en la entrevista es que la fórmula del éxito de una época puede convertirse en el veneno de la siguiente. Para una empresa de modelos, la inercia estratégica y los mecanismos de recompensa de la organización pueden ser más peligrosos que una carencia técnica puntual.

\subsection{Gemini, Meta y xAI}

Se valora que Gemini tiene modelos de gran capacidad y un buen nicho en el ecosistema, pero va muy rezagado en Coding; el error estratégico de Google fue no haber hecho de Coding su línea principal antes. Meta se considera un nuevo retador que sustituye a xAI como cuarto cabeza de serie; xAI, por su parte, se considera rezagada a corto plazo por sus vaivenes estratégicos y la salida de miembros de su founding team. Esto refleja que la competencia por los modelos de frontera ya no depende solo de los parámetros del modelo, sino de la organización, la estrategia, el producto y la capacidad de ejecución.

\begin{importantbox}{El verdadero problema de las tres grandes de Silicon Valley}
Una empresa de modelos líder que no dé importancia a Coding puede quedar fuera del primer nivel, porque Coding es la interfaz de ejecución con la que el modelo entra en el mundo digital y también la herramienta central para acelerar su propia investigación.
\end{importantbox}

\subsection{Resumen de la sección}

Las ventanas de las empresas de modelos rotan con rapidez. Quien logre convertir Coding/Agent en la línea principal de su organización podrá aprovechar el segundo acto; quien quede atrapado por el éxito de la etapa anterior podrá quedarse atrás.
\section{Harness Engineering: cómo entra el modelo en un flujo de trabajo real}

En la entrevista se menciona Harness Engineering: además del modelo, hace falta una capa de ingeniería que conecte el modelo con el flujo de trabajo real. Sin harness, el modelo es solo una interfaz potente; con harness, el modelo puede conectarse con el repositorio de código, las pruebas, los permisos, los registros, los reintentos, el despliegue y la retroalimentación.

\begin{figure}[H]
\centering
\includegraphics[width=0.92\textwidth]{harness-engineering.png}
\caption{Harness Engineering: la capa intermedia entre el modelo y el flujo de trabajo real.}
\end{figure}

\begin{knowledgebox}{Lectura de la figura: por qué no basta con que el modelo sea potente}
En la figura, el modelo es solo la primera capa. El harness se encarga de los prompts, las herramientas, los permisos, el estado y los reintentos; Workflow es el repositorio de código real y los procesos de negocio; Feedback devuelve al sistema las pruebas, la retroalimentación de los usuarios, los costos y los casos de fallo; solo entonces Improvement puede pasar al posentrenamiento y a la iteración del producto.
\end{knowledgebox}

\subsection{Términos clave: vocabulario relacionado con harness}

\noindent\begin{tabularx}{\textwidth}{p{0.22\textwidth}p{0.32\textwidth}X}
\toprule
Término & Problema que resuelve & Significado en este episodio \\
\midrule
Harness & Sistema externo que conecta el modelo con herramientas y flujos de trabajo & Determina si el modelo puede trabajar de forma estable. \\
Workflow & Flujo de tareas real, como issues, pruebas y despliegue & El valor del modelo se materializa, en última instancia, dentro del workflow. \\
Feedback & Pruebas, usuarios, costos, casos de fallo & Permite que el sistema aprenda y se mejore. \\
Permission & Permisos y límites de seguridad & Evita que el modelo realice operaciones sin autorización. \\
Retry / Recovery & Reintento y recuperación tras un error & Determina si las tareas de largo alcance pueden completarse de forma estable. \\
\bottomrule
\end{tabularx}

\subsection{Por qué se relaciona con el posentrenamiento}

El harness no es solo ingeniería de producto; también influye de vuelta en el entrenamiento. Un buen harness produce casos de fallo más claros, trayectorias de tareas más realistas y retroalimentación más verificable. Todo esto puede convertirse en datos de posentrenamiento. En el EP138, cuando Luo Fuli (罗福莉) habla del Agent post-train, enfatiza el entorno de tareas y el rollout; es la misma lógica que aquí.

\begin{warningbox}{No tratar el harness como un envoltorio de UI}
La UI es solo una pequeña parte del harness. Un harness real incluye gestión de estado, ejecución de herramientas, seguridad, registros, evaluación y retorno de la retroalimentación. Hacer solo una interfaz atractiva no permite que el modelo entre de forma confiable en producción.
\end{warningbox}

\subsection{Resumen de la sección}

Harness Engineering es el puente por el que el modelo se convierte en un sistema productivo. Conecta las capacidades del modelo con tareas reales y devuelve la retroalimentación de esas tareas al entrenamiento y a la iteración del producto.
\section{El modelo como nuevo sistema operativo}

La entrevista plantea que «el modelo es el sistema operativo de nueva generación». Esto no significa que el modelo sustituya a Linux, Windows o iOS, sino que el modelo podría convertirse en una nueva capa de orquestación entre la intención del usuario y la ejecución de las aplicaciones. Antes, el OS administraba el hardware, los archivos, los procesos y las aplicaciones; en el futuro, el modelo/Agent OS podría administrar intenciones, herramientas, contexto, permisos y el estado de las tareas.

\begin{figure}[H]
\centering
\includegraphics[width=0.92\textwidth]{model-as-os.png}
\caption{El modelo como OS de nueva generación: comprensión de la intención, orquestación de herramientas y memoria del estado.}
\end{figure}

\begin{knowledgebox}{Lectura de la figura: qué administra el Model OS}
En la figura, el OS del modelo ocupa el centro y conecta la intención del usuario, los servicios de las aplicaciones, la memoria de datos, el entorno de ejecución y los permisos de seguridad. La tesis que respalda es esta: el modelo no es solo una app, sino que podría convertirse en el nuevo punto de entrada a las aplicaciones y en la capa de orquestación de tareas.
\end{knowledgebox}

\subsection{Por qué Coding es la puerta de entrada para que el modelo funcione como OS}

El código es el lenguaje formal más potente del mundo digital. Si el modelo puede escribir código, modificarlo, llamar a API y ejecutar scripts, puede rebasar los límites de una sola app y convertirse en una capa intermedia que opera varios sistemas. Por eso, Coding es la puerta de entrada decisiva para que el modelo funcione como OS.

\begin{importantbox}{El cambio de app a OS}
Una app es un punto de entrada a una sola función; un OS es una capa de orquestación de recursos y tareas. Si el modelo puede comprender intenciones, llamar a herramientas, administrar el contexto y ejecutar tareas, pasará de la posición de app a la de OS.
\end{importantbox}

\subsection{Los tres grandes de China y el OS del modelo}

La entrevista menciona a los tres grandes de China, pero lo importante no es la clasificación concreta, sino una tesis estructural: si los equipos de modelos del país quieren aprovechar el segundo acto, necesitan comprender a la vez Coding, Agent, costos y escenarios de aplicación. La competencia por el OS del modelo no consiste solo en hacer un producto de chat, sino en integrarse con el trabajo de oficina, el desarrollo, el contenido, los datos, la automatización y los procesos empresariales.

\subsection{Resumen de la sección}

Que el modelo se convierta en el OS de nueva generación significa que pasa de ser un sistema que responde a ser un sistema que ejecuta tareas. Coding y el harness son la vía decisiva para que llegue a esa posición.
\section{Impacto social: deflación del trabajo de cuello blanco y ventana de desempleo}

La otra mitad del tono de este episodio proviene del impacto social. Los modelos de Coding multiplican por diez, o incluso por varias decenas, la productividad de los mejores investigadores e ingenieros; los ciclos de investigación e ingeniería se comprimen y las organizaciones, de manera natural, recalculan sus necesidades de personal. La deflación del trabajo de cuello blanco y la ventana de desempleo no son ciencia ficción lejana, sino un problema real: la capacidad de absorción de las organizaciones no alcanza a seguir el salto de productividad.

\begin{figure}[H]
\centering
\includegraphics[width=0.92\textwidth]{social-impact-chain.png}
\caption{Cómo el acelerador de Coding se transmite a la deflación del trabajo de cuello blanco y a la presión sobre el empleo.}
\end{figure}

\begin{knowledgebox}{Lectura de la figura: la cadena del impacto social}
La cadena comienza con el salto de los modelos de Coding, se transmite a la compresión de los ciclos de investigación y desarrollo, luego al cambio en la demanda de las organizaciones y, por último, se convierte en presión sobre los salarios y los puestos de trabajo. Lo importante no es la frase vacía «¿la IA reemplazará a las personas?», sino si la velocidad de difusión de la tecnología supera la velocidad con la que las organizaciones y las instituciones sociales pueden absorberla.
\end{knowledgebox}

\subsection{Por qué el trabajo de cuello blanco lo resiente primero}

Muchas tareas de cuello blanco, como Coding, redacción, análisis, producto, operaciones e investigación, ocurren en el mundo digital: sus entradas y salidas pueden expresarse como texto y código, por lo que es más fácil conectarlas a los modelos. Las tareas manuales y del mundo físico están limitadas por la robótica, el hardware y la complejidad del entorno, así que su difusión puede ser más lenta.

\begin{warningbox}{El aumento de productividad no se traduce automáticamente en mayor bienestar}
Si las ganancias de productividad se concentran en unas pocas empresas, en unos cuantos individuos sobresalientes o en el lado del capital, los trabajadores de cuello blanco comunes podrían experimentar primero presión salarial y reducción de puestos. Cómo se reparten los beneficios de la tecnología es un problema social, no algo que la capacidad de los modelos resuelva por sí sola.
\end{warningbox}

\subsection{Nuevas reflexiones sobre inversión}

En el plano de la inversión, este episodio contiene un juicio implícito: si Coding es un acelerador, las empresas de modelos y de cadenas de herramientas se revalorizarán; si el modelo se convierte en el OS, también se reorganizarán los puntos de entrada de las aplicaciones y el panorama de plataformas. Sin embargo, invertir no puede basarse solo en palabras de moda: hay que ver si el modelo entra en flujos de trabajo reales, si genera retroalimentación sostenible y si tiene ventajas de costo y canales de distribución.

\subsection{Resumen de la sección}

El impacto social de Coding/Agent proviene de la compresión de la productividad. Traerá nuevas empresas, nuevos puntos de entrada y nuevas oportunidades de inversión, pero también una revaloración de los puestos de cuello blanco, presión salarial y desafíos institucionales.
\section{Cronología del trimestre: por qué este trimestre no es una actualización ordinaria}

Este informe trimestral describe el primer trimestre de 2026 como una etapa con una sensación de aceleración muy intensa, como la que empuja al pasajero contra el asiento: los modelos pasan del chat a las tareas agentic, Coding comprime los ciclos de investigación y desarrollo, y la frecuencia con la que los laboratorios de frontera lanzan productos y features aumenta notablemente. Lo decisivo aquí no es un modelo aislado, sino que las capacidades, las herramientas, las organizaciones y los flujos de trabajo cambian al mismo tiempo.

\subsection{El hilo conductor de los informes trimestrales de 2023 a 2026}

\noindent\begin{longtable}{p{0.20\textwidth}p{0.36\textwidth}p{0.35\textwidth}}
\toprule
Etapa & Hilo conductor & Cómo lo retoma este episodio \\
\midrule
2023 & ChatGPT inaugura el paradigma del chat; el código abierto empieza a alcanzarlo & El primer acto es «el modelo sabe hablar y sabe responder». \\
2024 & Surgen la infraestructura de AGI, el razonamiento, el self-play RL y el relato de o1 & Los modelos empiezan a pasar de responder a razonar y actuar. \\
2025 & Diferenciación y convergencia, suites completas de productos, integración vertical, AI War & Las empresas de modelos empiezan a integrarse en torno a productos y ecosistemas. \\
2026 Q1 & Coding/Agent se convierte en el hilo conductor del segundo acto & Los modelos entran en los flujos de trabajo digitales y empiezan a transformar las organizaciones de investigación y desarrollo. \\
\bottomrule
\end{longtable}

\begin{importantbox}{El valor de un informe trimestral}
Una noticia aislada tiende a exagerar el ruido de corto plazo; la perspectiva trimestral permite observar qué juicios se sostienen en el tiempo. El juicio de continuidad más importante de este episodio es este: el año pasado Coding todavía era una predicción; este año ya se convirtió en un tsunami.
\end{importantbox}

\subsection{Resumen de la sección}

La importancia de EP136 consiste en enlazar los juicios de los trimestres anteriores: los modelos pasaron del chat al razonamiento y a los ecosistemas de productos, hasta llegar a Coding/Agent, una etapa capaz de cambiar directamente los flujos de trabajo.
\section{Por qué Coding se parece a una GPU: un acelerador, no una aplicación}

Comparar el modelo de Coding con una GPU es la analogía con mayor poder explicativo de este episodio. La GPU no es en sí una aplicación de consumo, pero hace posibles el entrenamiento y la inferencia; del mismo modo, el modelo de Coding no es solo un complemento del IDE, sino una infraestructura que acelera la investigación, la ingeniería, el procesamiento de datos, los experimentos y la iteración de productos.

\subsection{Tres vías de aceleración}

\noindent\begin{tabularx}{\textwidth}{p{0.22\textwidth}p{0.34\textwidth}X}
\toprule
Vía de aceleración & Dónde ocurre & Efecto \\
\midrule
Amplificación de la productividad individual & Investigadores, programadores y arquitectos de primer nivel & El 1\% del talento se amplifica de 10 a 50 veces y se acorta el ciclo que va de la idea al experimento. \\
Capacidad de producción de I+D de la organización & feature, pipeline de datos, experimentos multimodales & Iteraciones que antes tomaban semanas se reducen a días y aumenta la frecuencia de lanzamiento de productos. \\
IA que acelera la IA & El modelo ayuda a escribir código de entrenamiento, scripts de evaluación y a depurar experimentos & La propia investigación en IA se acelera con herramientas de IA, lo que forma un ciclo de retroalimentación recursivo. \\
\bottomrule
\end{tabularx}

\begin{knowledgebox}{Por qué el código es la palanca del mundo digital}
El lenguaje natural expresa problemas; el código expresa soluciones. Siempre que una tarea ocurra dentro de un sistema digital, el código puede conectar datos, herramientas, entornos, permisos y ejecución. Que un modelo domine el coding equivale a que tenga en sus manos la palanca de mando del mundo digital.
\end{knowledgebox}

\subsection{Resumen de la sección}

Coding es un acelerador, no una aplicación aislada. Amplifica a las personas, a las organizaciones y a la propia investigación en IA; por eso podría convertirse en la infraestructura fundamental del segundo acto de la AGI.

\section{La dependencia de la trayectoria de los tres grandes de Silicon Valley}

Lo central en la valoración que este episodio hace de Anthropic, OpenAI y Gemini no es el chisme, sino la dependencia de la trayectoria. El éxito de cada empresa en la etapa anterior moldea la atención de su organización en la etapa siguiente. La cultura de underdog de Anthropic y su inversión en coding/agent la colocaron frente a la ventana de oportunidad; el éxito de ChatGPT en ToC podría llevar a OpenAI a descuidar coding; Gemini tiene un ecosistema y una capacidad de modelos sólidos, pero en lo estratégico se parece más a un seguidor aventajado.

\subsection{Comparación punto por punto}

\noindent\begin{longtable}{p{0.18\textwidth}p{0.25\textwidth}p{0.25\textwidth}p{0.22\textwidth}}
\toprule
Empresa & Fortalezas según la entrevista & Riesgos según la entrevista & Lección para el segundo acto de Coding \\
\midrule
Anthropic & Atención al detalle, cultura exigente, ventana fuerte en Agent/Coding & Su ventaja puede ser alcanzada; escalar sigue siendo un reto & Invertir pronto en flujos de trabajo abre una ventana de oportunidad. \\
OpenAI & Marca, modelos, producto y capital sólidos & El éxito en ToC puede convertirse en inercia estratégica & El éxito pasado puede ocultar la nueva línea principal. \\
Gemini & Ecosistema, recursos y capacidad de modelos sólidos & Rezago en Coding, poco impulso & El ecosistema no sustituye a una estrategia clara. \\
Meta & Código abierto, capital, talento y posición de retadora & El ciclo cerrado de producto es inestable & El código abierto puede ser una vía para que un rezagado alcance a los demás. \\
xAI & Muchos recursos y mucha atención & Vaivenes estratégicos, salida de personal del equipo & La estabilidad de la organización importa más que el relato de las estrellas. \\
\bottomrule
\end{longtable}

\begin{warningbox}{No tomes los juicios sobre las empresas como conclusiones permanentes}
Este episodio es una instantánea trimestral, no una clasificación a largo plazo. La ventana de oportunidad de las empresas de modelos puede cambiar de un mes a otro. Un método de análisis más estable consiste en observar la dependencia de la trayectoria: qué recompensa la organización, qué descuida y quién puede adoptar más rápido un nuevo paradigma.
\end{warningbox}

\subsection{Resumen de la sección}

En el fondo, comparar a los tres grandes de Silicon Valley es comparar su atención estratégica. El segundo acto de Coding recompensará a las empresas que pongan el modelo, el producto, los datos y el ciclo cerrado de ingeniería sobre una misma línea principal.
\section{Los tres grandes de China y la ventaja en costos}

En este episodio también se menciona a los tres grandes de China. Aunque la transcripción no entra en muchos detalles, por el contexto anterior y posterior la ventaja de las empresas chinas de modelos parece venir más bien de los costos, la velocidad de ingeniería, los escenarios de aplicación y el ecosistema de desarrolladores, y no de aplastar a la competencia en un solo benchmark. Los escenarios de Coding/Agent son especialmente sensibles al costo, porque las tareas de largo alcance multiplican cada llamada al modelo.

\subsection{Por qué el costo es una variable estratégica}

En un escenario de chat, quizá el usuario todavía acepte un modelo algo más caro; en un escenario de Agent, una sola tarea puede requerir decenas de rondas de llamadas, ejecución de herramientas y verificación. Si el costo por llamada es demasiado alto, al producto le cuesta escalar. Si los modelos nacionales baratos y eficaces logran conectarse a un buen harness, podrían obtener una ventaja en un gran volumen de tareas de frecuencia media y alta.

\begin{importantbox}{La ventaja en costos no es una ventaja de gama baja}
Un modelo barato no está limitado a tareas de poco valor. Un framework de Agent puede usar el modelo barato en los pasos de alta frecuencia y el modelo potente en las decisiones críticas. El enrutamiento entre varios modelos convierte la estructura de costos en un factor de competitividad.
\end{importantbox}

\subsection{Resumen de la sección}

La oportunidad de las empresas chinas de modelos consiste en lo siguiente: si logran incorporar su ventaja en costos, su velocidad de ingeniería y la retroalimentación de sus aplicaciones a los flujos de trabajo de Coding/Agent, podrían ocupar una posición diferenciada en el segundo acto.
\section{Respuesta social: qué hacer después de la deflación del trabajo de oficina}

La parte más sombría de este episodio es la deflación del trabajo de oficina y el periodo de desempleo que abre. En lo técnico, los modelos de Coding acortan los ciclos de desarrollo; en lo organizacional, unas pocas personas pueden completar más tareas; en lo económico, el precio del trabajo de oficina se volverá a fijar. La pregunta no es «¿reemplazará la IA a las personas?», sino si la sociedad puede ofrecer un nuevo lugar a quienes sean reemplazados o reorganizados.

\subsection{Tres grupos que recibirán el impacto primero}

\noindent\begin{tabularx}{\textwidth}{p{0.22\textwidth}p{0.34\textwidth}X}
\toprule
Grupo & Por qué lo afecta primero & Dirección de respuesta \\
\midrule
Trabajadores del conocimiento con tareas repetitivas & Su trabajo puede expresarse en texto, tablas y código & Pasar a la definición de tareas, la revisión, la comprensión del cliente y el diseño de procesos. \\
Programadores principiantes & El modelo puede generar gran parte del código de implementación & Aprender arquitectura, pruebas, review y comprensión de sistemas. \\
Coordinadores de mando medio & Las herramientas pueden automatizar el reenvío de información y el seguimiento del avance & Pasar al juicio, la asignación de recursos y la resolución de conflictos entre equipos. \\
\bottomrule
\end{tabularx}

\begin{knowledgebox}{El verdadero objetivo de la recapacitación}
La recapacitación no debe limitarse a enseñar a «usar herramientas de IA»; debe formar a las personas para definir problemas, verificar resultados, organizar flujos de trabajo y comprender el negocio. Las herramientas cambiarán; el criterio y la estructura de responsabilidades cambian más despacio.
\end{knowledgebox}

\subsection{Resumen de la sección}

La deflación del trabajo de oficina resulta de que la tecnología se difunde más rápido de lo que las organizaciones pueden absorberla. La respuesta social no puede depender solo del aprendizaje individual: también hace falta que la organización de las empresas y el diseño institucional redistribuyan los beneficios del aumento de productividad.
\section{La estructura de producto del modelo como OS: de la puerta de entrada del chat a la puerta de entrada de las tareas}

Llamar al modelo una nueva generación de OS es lo que más fácilmente se malinterpreta como una exageración. Dicho con más precisión, el modelo podría convertirse en la «puerta de entrada de las tareas»: el usuario ya no abre primero una app, sino que primero expresa un objetivo, y el modelo elige las herramientas, invoca los servicios, organiza el contexto, ejecuta las acciones y devuelve el resultado. Este cambio reconfigurará la distribución de aplicaciones, la forma del SaaS y la lógica de compra del software empresarial.

\subsection{Estructura de producto en tres capas}

\noindent\begin{longtable}{p{0.22\textwidth}p{0.34\textwidth}p{0.35\textwidth}}
\toprule
Capa & Responsabilidad & Foco de la competencia \\
\midrule
Capa del modelo & Comprender la intención, planificar tareas, generar código o acciones & Capacidad base, capacidad de Coding, capacidad de uso de herramientas. \\
Capa de Harness & Conectar aplicaciones, permisos, memoria, registros, reintentos, evaluación & Confiabilidad del flujo de trabajo, costo, seguridad y observabilidad. \\
Capa de aplicación & Alojar los objetos de negocio concretos y la relación con el usuario & Datos, escenarios, distribución, hábitos del usuario y pago. \\
\bottomrule
\end{longtable}

\begin{knowledgebox}{Por qué la conversión en OS amenaza la puerta de entrada de las aplicaciones}
Si el usuario expresa sus objetivos a través del modelo, y el modelo se encarga de elegir las aplicaciones e invocar los servicios, el valor de las apps tradicionales como puerta de entrada disminuirá. Las aplicaciones siguen siendo importantes, pero podrían pasar de «el usuario las abre por iniciativa propia» a «el modelo las asigna según la tarea».
\end{knowledgebox}

\subsection{Los cambios en el software empresarial}

El software empresarial se diseñaba en torno a interfaces operadas por personas: formularios, botones, flujos, permisos e informes. En la era de los Agent, el software empresarial también necesita ofrecer API que el modelo pueda invocar, un estado legible, registros con capacidad de auditoría y operaciones recuperables. Quien convierta antes su producto en un agent-friendly workflow tendrá más facilidad para entrar en la red de asignación de tareas del modelo como OS.

\begin{importantbox}{Características del software agent-friendly}
API claras, permisos explícitos, estado verificable, acciones que se pueden revertir, registros estructurados, sandbox de bajo costo y buena documentación. Sin esto, por muy capaz que sea el modelo, difícilmente podrá operar los sistemas empresariales de forma segura y estable.
\end{importantbox}

\subsection{Resumen de la sección}

El modelo como OS no sustituye al núcleo del sistema operativo, sino que se convierte en la capa de asignación de tareas. La competencia futura del software pasará de «quién controla la puerta de entrada de la interfaz» a «quién puede ser invocado de forma confiable por el modelo».
\section{Marco de inversión: cómo valorar una empresa de Coding/Agent}

Al final del episodio se abordan nuevas reflexiones sobre inversión. Si Coding es el acelerador de la IA, la decisión de inversión no puede basarse solo en el crecimiento de usuarios o en las tablas de clasificación de modelos; hay que ver si una empresa controla el flujo de trabajo, la retroalimentación y la estructura de costos. Si una empresa de Coding/Agent solo es un envoltorio que llama a un modelo, su barrera de entrada es muy débil; si logra acumular datos de tareas reales, casos de fallo, integraciones de herramientas y flujos de trabajo de los usuarios, su barrera de entrada será más sólida.

\subsection{Cuatro dimensiones de evaluación}

\noindent\begin{tabularx}{\textwidth}{p{0.23\textwidth}p{0.34\textwidth}X}
\toprule
Dimensión & Pregunta que hay que hacer & Señal favorable \\
\midrule
Profundidad en el flujo de trabajo & ¿Forma parte de las tareas que el usuario tiene que completar todos los días? & Integración profunda con el repositorio de código, las pruebas, el despliegue y los flujos de datos. \\
Ciclo de retroalimentación & ¿Acumula trayectorias de fallos y de éxitos que sirvan para entrenar? & Cuenta con registros estructurados, evaluaciones y datos de correcciones hechas por personas. \\
Estructura de costos & ¿Puede reducir el costo por tarea con enrutamiento entre varios modelos? & El modelo grande solo se ocupa de los pasos críticos y los modelos pequeños atienden las tareas de alta frecuencia. \\
Canal de distribución & ¿Tiene hábitos de uso o un nicho propio en el ecosistema? & IDE, sistemas empresariales, navegador, puntos de acceso para la colaboración en equipo. \\
\bottomrule
\end{tabularx}

\begin{warningbox}{Trampas comunes en el discurso de inversión}
«Somos una empresa de Agent» no es una barrera de entrada; «integramos el modelo más reciente» tampoco lo es. La verdadera barrera de entrada suele estar en la profundidad en el flujo de trabajo, la retroalimentación de datos, la capacidad de ejecución de la organización y el canal de distribución, no en el envoltorio del prompt.
\end{warningbox}

\subsection{Resumen de la sección}

La inversión en Coding/Agent tiene que pasar de la demo al workflow. Que una empresa logre acumular retroalimentación, reducir costos y ocupar el punto de acceso importa más que las demostraciones de capacidad de los modelos a corto plazo.
\section{Impacto en el empleo por niveles: qué puestos cambian primero}

La deflación del trabajo de cuello blanco no ocurre de manera uniforme. Coding/Agent afectará primero a los puestos cuyas tareas tienen un alto grado de digitalización, cuyos resultados son verificables y cuyos procesos pueden descomponerse. No necesariamente eliminará de inmediato una profesión completa, pero primero reorganizará la proporción de tareas dentro de cada puesto.

\subsection{Matriz de impacto por puesto}

\noindent\begin{longtable}{p{0.22\textwidth}p{0.34\textwidth}p{0.35\textwidth}}
\toprule
Puesto/tarea & Parte que se comprime primero & Parte más difícil de sustituir \\
\midrule
Programador júnior & Código repetitivo, bugs simples, scripts, encapsulamiento de interfaces & Criterio de arquitectura, comprensión de sistemas complejos, responsabilidad por la puesta en producción. \\
Análisis de datos & Limpieza, SQL, reportes, primeras versiones de visualizaciones & Definición de métricas, interpretación del negocio, juicio causal. \\
Producto/operaciones & Documentación, resúmenes de la competencia, configuración de procesos & Comprensión profunda de los usuarios, toma de decisiones entre alternativas, coordinación organizacional. \\
Asistente de investigación & Organización de bibliografía, scripts de experimentos, síntesis de resultados & Planteamiento de preguntas, diseño experimental, explicación de los fracasos. \\
Gestión y coordinación & Sincronización del estado, minutas de reuniones, calendarización & Resolución de conflictos, decisiones sobre recursos, asunción de responsabilidades. \\
\bottomrule
\end{longtable}

\begin{importantbox}{La forma real del cambio en las profesiones}
La IA no necesariamente sustituye un puesto de una sola vez, sino que primero sustituye tareas. Los puestos se reorganizarán: las partes de bajo criterio, alta repetición y resultados verificables se automatizan; las partes de alta responsabilidad, mucho contexto y decisiones difíciles entre alternativas cobran más importancia.
\end{importantbox}

\subsection{Resumen de la sección}

El impacto en el empleo debe analizarse con la granularidad de las tareas. Lo primero que se comprime son las tareas digitalizadas, repetitivas y verificables; lo más difícil de sustituir son la responsabilidad, el criterio, la organización y las decisiones sobre valores.
\section{Conexión con EP139, EP138 y EP137}

EP139 trata la historia técnica de los Agent y explica por qué el digital agent se convertirá en un nuevo paradigma; EP138 trata cómo los laboratorios de modelos reorganizan el posentrenamiento, la capacidad de cómputo y la organización en función de los Agent; EP137 trata cómo AI for Math lleva el ciclo cerrado al estilo de los Agent hacia el descubrimiento científico; EP136, por su parte, sitúa estos juicios en el panorama trimestral de la industria y se pregunta cómo se reordenarán las empresas de modelos, las aplicaciones, el empleo y la inversión.

\begin{knowledgebox}{Ruta de lectura de los cuatro episodios}
Primero lee EP139 para construir el marco técnico de los Agent; después lee EP138 para entender cómo actúan los laboratorios de modelos; luego lee EP137 para ver cómo la IA entra en la investigación matemática; por último, lee EP136 para devolver estos cambios a los planos de la industria, las empresas y la sociedad.
\end{knowledgebox}

\subsection{Resumen de la sección}

EP136 es un panorama industrial de los episodios anteriores: traduce los cambios técnicos de los Agent, el posentrenamiento y AI for Math en estrategia empresarial, OS de modelos e impacto social.
\section{Términos clave: índice de conceptos de este episodio}

\begin{longtable}{p{0.24\textwidth}p{0.32\textwidth}p{0.35\textwidth}}
\toprule
Término & Explicación breve & Función en este episodio \\
\midrule
Modelo de Coding & Modelo capaz de entender, generar, modificar y depurar código & Se considera el acelerador del segundo acto de la AGI. \\
Agentic Workflow & El modelo completa una tarea en varios pasos dentro de herramientas y entornos & Hace que el modelo pase de conversar a realizar trabajo. \\
Harness Engineering & Capa de ingeniería que conecta el modelo con flujos de trabajo reales & Determina si la capacidad del modelo puede llevarse a la práctica de forma estable. \\
Model OS & El modelo como capa de orquestación que va de la intención a la ejecución de aplicaciones & Explica por qué el modelo podría redefinir el punto de entrada a las aplicaciones. \\
Deflación del trabajo de cuello blanco & La automatización presiona la oferta, la demanda y los salarios del trabajo de cuello blanco & Eje principal del impacto social. \\
Forward Deployment & Ingenieros que despliegan sistemas de IA trabajando de cerca con el cliente & Muestra que llevar los Agent a la práctica todavía requiere un servicio intensivo. \\
Coding as GPU & El modelo de Coding acelera la investigación y la ingeniería como lo hace una GPU & La metáfora con mayor poder explicativo de este episodio. \\
\bottomrule
\end{longtable}

\subsection{Resumen de la sección}

Los términos de este episodio giran en torno a una idea central: el Coding convierte al modelo en un acelerador de la productividad e impulsa una nueva valoración de las empresas de modelos, del ecosistema de aplicaciones y de la estructura social.
\section{Apéndice: corrección de la transcripción y ruta de repaso}

Este episodio no tiene subtítulos de la plataforma, y la transcripción automática confundió varios nombres de modelos y términos en inglés. Al preparar la nota hay que unificar las aproximaciones fonéticas de la transcripción con los términos estándar; de lo contrario, se dificultan las búsquedas posteriores y la comparación entre episodios.

\subsection{Términos clave: corrección de confusiones frecuentes}

\noindent\begin{longtable}{p{0.24\textwidth}p{0.31\textwidth}p{0.36\textwidth}}
\toprule
Término estándar & Confusión frecuente / escritura aproximada & Motivo de la corrección \\
\midrule
Chatbot & Charbot, 拆的 GTT & Forma de interacción del primer acto. \\
Anthropic & AnswerPick, Ansropic & Una de las tres grandes empresas de Silicon Valley; empresa central de este episodio. \\
Claude Code & Cloud Code, Cowl Code & Herramienta representativa de los flujos de trabajo de Coding/Agent. \\
Codex & CodeX, denominación genérica de los modelos de código & Ecosistema de herramientas de coding de OpenAI. \\
Gemini & Germline, Gemline & Ecosistema de modelos de frontera de Google. \\
Harness Engineering & Harness, 哈尼斯工程 & Capa intermedia por la que el modelo entra en los flujos de trabajo reales. \\
Model as OS & sistema operativo de modelo, modelo OS & Juicio central de este episodio sobre la forma futura de los modelos. \\
White-collar deflation & deflación de cuello blanco & Palabra clave de la parte sobre el impacto social. \\
\bottomrule
\end{longtable}

\subsection{Ruta de repaso}

Si solo se lee una vez, se recomienda repasar en cuatro pasos. Primero, leer las secciones 2--3 para entender por qué Coding pasó de ser una aplicación a ser un acelerador. Segundo, leer las secciones 4--6 para entender la dependencia de trayectoria de las empresas de modelos, el harness y el model OS. Tercero, leer las secciones 7--10 para entender el impacto social, el marco de inversión y la reorganización de los puestos de trabajo. Cuarto, conectar este episodio con EP139, EP138 y EP137 para ver la relación entre la historia técnica de los Agent, los laboratorios de modelos, AI for Math y los informes trimestrales de la industria.

\begin{knowledgebox}{El resumen más breve de este episodio}
Coding es el lenguaje de ejecución del mundo digital; Agent es la forma en que el modelo entra en el entorno; Harness es el puente por el que el modelo entra en los flujos de trabajo; Model OS es la visión final en la que el modelo se convierte en el punto de entrada de las tareas; la deflación de cuello blanco es la presión social que sigue al salto de productividad.
\end{knowledgebox}

\subsection{Resumen de la sección}

La corrección de la transcripción estabiliza los términos, y la ruta de repaso permite reutilizar los juicios trimestrales. Al seguir preparando la serie de Zhang Xiaojun AI, este episodio puede servir como plantilla del tipo «informe trimestral de la industria/reseña».
\section{Síntesis y ampliación}

\subsection{Conclusiones centrales}

\begin{enumerate}
\item Coding es el segundo acto de la AGI, porque lleva al modelo de conversar a ejecutar tareas en el mundo digital.
\item La ventana de competencia entre las empresas de modelos de frontera se está reordenando; el éxito de la etapa anterior puede convertirse en dependencia de la trayectoria en la siguiente.
\item La Harness Engineering determina si el modelo puede entrar en flujos de trabajo reales.
\item El modelo se vuelve la pieza central de una nueva generación de sistemas operativos porque puede coordinar herramientas, contexto, permisos y estado de las tareas.
\item Coding/Agent traerá una deflación del trabajo de oficina y una ventana de desempleo; la sociedad podría absorber el cambio más despacio de lo que se difunde la tecnología.
\end{enumerate}

\subsection{Preguntas abiertas}

\begin{itemize}
\item ¿Llegarán los modelos de Coding a ser, como la GPU, el acelerador básico de toda la investigación en IA?
\item ¿El sistema operativo basado en modelos quedará en manos de las grandes empresas actuales o de una nueva generación de empresas de aplicaciones AI-native?
\item ¿En qué industrias se producirá primero la deflación del trabajo de oficina y cómo redistribuirá la sociedad los beneficios de la productividad?
\item ¿Podrán las empresas chinas de modelos alcanzar el segundo acto apoyándose en sus costos, sus aplicaciones y su ecosistema de desarrolladores?
\end{itemize}

\subsection{Lecturas adicionales}

\begin{itemize}
\item EP139, historia técnica de los Agent: para entender por qué Coding/Agent forma parte de una genealogía tecnológica más amplia.
\item EP138, entrevista con Luo Fuli (罗福莉): para entender cómo el paradigma Agent cambia el posentrenamiento y la asignación de capacidad de cómputo.
\item Práctica con herramientas como Claude Code, Codex y OpenClaw: para observar cómo los modelos de Coding cambian el flujo de trabajo de investigación y desarrollo.
\item Harness engineering, agent evaluation, computer use agent benchmark: para entender la capa intermedia por la que los modelos entran en los sistemas de producción.
\end{itemize}

\begin{importantbox}{Juicio final}
Si en 2023 la pregunta era «¿el modelo puede hablar?», en 2026 la pregunta es «¿el modelo puede trabajar?». Coding es el primer terreno firme de esta transición; quien logre conectarlo a flujos de trabajo reales estará más cerca de la próxima generación de sistemas operativos de IA.
\end{importantbox}

\end{document}
